
\chapterwithtoc{Résumé}
\label{chap:resume}

Dans un centre d'opérations de sécurité, le goulot d'étranglement n'est plus la
détection mais la qualification de ce qui est détecté. Rattacher une activité
observée au référentiel MITRE ATT\&CK reste un travail manuel, lent et
tributaire d'une expertise rare, et la difficulté redouble chez un exploitant
d'infrastructure énergétique, dont les incidents traversent la frontière entre
informatique de gestion et informatique industrielle.

Cet essai conçoit et mesure un écosystème agentique ancrant un modèle de langage
sur ATT\&CK au moyen du \emph{Model Context Protocol}. Un corpus de
42 échantillons est produit en laboratoire, à l'abri de la contamination qui
fausse les bancs d'essai publics, sur deux zones --- systèmes d'information et
automates --- et deux matrices, Enterprise et ICS. Dix modèles de la génération
courante y sont évalués sans ancrage, en trois exécutions par échantillon, soit
1\,260 appels, selon une méthodologie qui sépare cinq catégories de prédiction
et deux bases de calcul. Le modèle retenu est ensuite rejoué avec le serveur, à
corpus et conditions identiques.

Sans ancrage, aucun modèle ne dépasse 0,603 de F1 au niveau technique, et la
performance chute d'un tiers sur la matrice industrielle. Les identifiants
fabriqués sont exclusivement des tactiques, et exclusivement sur des journaux
industriels : l'erreur est un défaut de repérage entre univers, non une
invention arbitraire. Avec ancrage, le F1 technique passe de 0,563 à 0,706, la
matrice industrielle double presque, et le taux d'identifiants périmés ne
diminue pas : il disparaît. La contribution ne tient pas à l'accès à une source,
mais à quatre décisions de conception --- version épinglée et jointe à chaque
réponse, résolution des révocations, granularité d'outils bornée, trace d'usage
vérifiable.

\vspace{1em}
\noindent\textbf{Mots-clés :} MITRE ATT\&CK, Model Context Protocol, grands
modèles de langage, ancrage, centre d'opérations de sécurité, technologies
opérationnelles, ICS, évaluation.
